% !TeX program = xelatex
\ifdefined\BookMaster\else
\documentclass[11pt,oneside]{article}
% !TeX program = xelatex
\usepackage[b5paper,left=22mm,right=22mm,top=20mm,bottom=22mm,footskip=11mm]{geometry}
\usepackage{fontspec}
\setmainfont{TeX Gyre Termes}
\setsansfont{TeX Gyre Heros}
\usepackage{amsmath,amssymb,mathtools,amsthm}
\renewcommand{\cdots}{\mathinner{\cdotp\cdotp\cdotp}}
\usepackage{graphicx,adjustbox,array,booktabs,longtable,multirow,enumitem,changepage,needspace}
\usepackage[table]{xcolor}
\usepackage{fancyhdr}
\raggedbottom
\interfootnotelinepenalty=10000
\newcommand{\BookFootnote}[2]{\begingroup\renewcommand{\thefootnote}{\fnsymbol{footnote}}\footnote[#1]{#2}\endgroup}
\newcommand{\BookFootnoteMark}[1]{\begingroup\renewcommand{\thefootnote}{\fnsymbol{footnote}}\footnotemark[#1]\endgroup}
\newcommand{\BookFootnoteText}[2]{\begingroup\renewcommand{\thefootnote}{\fnsymbol{footnote}}\footnotetext[#1]{#2}\endgroup}
\newenvironment{BookFloat}[1]{\par\addvspace{.6\baselineskip}\noindent\def\BookFloatID{#1}\begin{minipage}{\linewidth}\phantomsection\label{bookfloat:#1:start}}{\phantomsection\label{bookfloat:\BookFloatID:end}\end{minipage}\par\addvspace{.6\baselineskip}}

\usepackage[unicode,hidelinks]{hyperref}
\pagestyle{fancy}
\fancyhf{}
\newcommand{\BookFooter}{}
\fancyfoot[L]{\parbox[b]{\dimexpr\textwidth-15mm}{\footnotesize\BookFooter}}
\fancyfoot[R]{\footnotesize\thepage}
\renewcommand{\headrulewidth}{0pt}
\renewcommand{\footrulewidth}{0pt}
\setlength{\parindent}{0pt}
\setlength{\parskip}{0.2em}
\setlength{\emergencystretch}{3em}
\setlist{nosep,leftmargin=1.7em,labelsep=.65em}
\newif\ifBookNarrative
\BookNarrativetrue
\newenvironment{BookBody}{\setlength{\parindent}{0pt}}{}
\newcommand{\BookStatement}[1]{\par\addvspace{.6\baselineskip}\noindent #1\enspace\ignorespaces}
\newcommand{\BookLabel}[1]{\par\addvspace{.6\baselineskip}\needspace{3\baselineskip}\noindent\textbf{#1}\quad\ignorespaces}
\newcommand{\BookHeading}[1]{\par\addvspace{\baselineskip}\needspace{4\baselineskip}\noindent{\large\bfseries #1}\par\medskip}
\newcommand{\BookSubheading}[1]{\par\addvspace{.7\baselineskip}\needspace{3\baselineskip}\noindent{\bfseries #1}\par\smallskip}
\newcommand{\BookTextHeading}[1]{\par\addvspace{.7\baselineskip}\needspace{3\baselineskip}\noindent #1\par\smallskip}
\newcommand{\BookSectionTitle}[2]{\par\addvspace{1.2\baselineskip}\needspace{5\baselineskip}\noindent\colorbox{gray!25}{\sffamily\bfseries\small SECTION #1}\quad{\large\bfseries #2}\par\medskip}
\newcommand{\BookExerciseTitle}[1]{\par\addvspace{\baselineskip}\needspace{4\baselineskip}\noindent{\color{gray}\rule{1.5mm}{1.5mm}}\enspace{\large\bfseries Exercises #1}\par\medskip}
\newcommand{\BookBullet}[1]{\begin{list}{\textbullet}{\setlength{\leftmargin}{1.7em}\setlength{\labelwidth}{1em}\setlength{\labelsep}{.7em}\setlength{\topsep}{0pt}\setlength{\parsep}{0pt}\setlength{\itemsep}{0pt}}\item #1\end{list}}
\newcommand{\BookExercise}[2]{\begin{list}{\textbf{#1}}{\setlength{\leftmargin}{2.2em}\setlength{\labelwidth}{1.5em}\setlength{\labelsep}{.7em}\setlength{\topsep}{.15em}\setlength{\parsep}{0pt}\setlength{\itemsep}{0pt}}\item #2\end{list}}
\newcommand{\BookEnd}[1]{\unskip\nobreak\hfill\ensuremath{#1}}
\newcommand{\BookSource}[2]{\phantomsection\label{sourcepage:#1}} % Source-page anchors, retained through reflow.
\usepackage{multicol}
\usepackage{tikz}
\newcommand{\BookPartEntry}[2]{\par\noindent\makebox[17mm][l]{\bfseries Section #1}#2\par}
\newcommand{\BookPartTitle}[3]{\par\bigskip\noindent\begin{tikzpicture}\node[anchor=north west,inner sep=0pt,outer sep=0pt] (bookpart) at (0,0) {\begin{minipage}{\dimexpr\linewidth-.5pt}\hspace*{26mm}\begin{minipage}[t]{\dimexpr\linewidth-26mm}\begin{adjustbox}{max width=\linewidth}{\fontsize{20}{24}\selectfont\bfseries #2}\end{adjustbox}\par\vspace{6mm}\hrule\vspace{6mm}{\sffamily\footnotesize\setlength{\parskip}{.3em}#3}\vspace{3mm}\end{minipage}\end{minipage}};\path[shade,top color=gray!60,bottom color=white] (0,0) rectangle ([xshift=24mm]bookpart.south west);\node[anchor=north,font=\sffamily\bfseries\small] at (12mm,-2mm) {PART};\node[anchor=north,font=\sffamily\bfseries\Huge] at (12mm,-7mm) {#1};\end{tikzpicture}\par\bigskip}
\newsavebox{\BookMathBox}
\makeatletter
\newcommand{\BookInline}[1]{\leavevmode\hskip0pt\relax\begingroup\in@{\begin}{#1}\ifin@\sbox{\BookMathBox}{\(#1\)}\ifdim\wd\BookMathBox>\dimexpr\linewidth-1em\relax\resizebox{\dimexpr\linewidth-1em\relax}{!}{\usebox{\BookMathBox}}\else\(#1\)\fi\else\in@{\left}{#1}\ifin@\sbox{\BookMathBox}{\(#1\)}\ifdim\wd\BookMathBox>\dimexpr\linewidth-1em\relax\resizebox{\dimexpr\linewidth-1em\relax}{!}{\usebox{\BookMathBox}}\else\(#1\)\fi\else\(#1\)\fi\fi\endgroup\hskip0pt\relax}
\makeatother
\DeclareMathOperator{\Aut}{Aut}
\DeclareMathOperator{\End}{End}
\DeclareMathOperator{\Hom}{Hom}
\DeclareMathOperator{\GF}{GF}
\DeclareMathOperator{\irr}{irr}
\DeclareMathOperator{\lt}{lt}
\DeclareMathOperator{\lp}{lp}
\providecommand{\textcircled}[1]{\textnormal{(#1)}}

\begin{document}
\fi
\renewcommand{\BookFooter}{Section 34 Unique Factorization Domains}
\pdfbookmark[0]{Section 34 Unique Factorization Domains}{34_unique_factorization_domains}
\ifdefined\BookMaster\else\setcounter{page}{275}\fi
% Source PDF page 292; original printed page 275
\BookSource{292}{275}
\BookNarrativetrue

\begin{BookBody}

\BookSectionTitle{34}{Unique Factorization Domains}

\noindent The integral domain \BookInline{\mathbb{Z}} is our standard example of an integral domain in which there is unique factorization into primes (irreducibles). Section 28 showed that for a field \BookInline{F,\allowbreak F[x]} is also such an integral domain with unique factorization. In order to discuss analogous 
% Source PDF page 293; original printed page 276
\BookSource{293}{276}
ideas in an arbitrary integral domain, we shall give several definitions, some of which are repetitions of earlier ones. It is nice to have them all in one place for reference.

\end{BookBody}

\BookNarrativetrue

\begin{BookBody}

\BookLabel{34.1 Definition} Let \BookInline{R} be a commutative ring with unity and let \BookInline{a,\allowbreak b\in R.} If there exists \BookInline{c\in R} such that \BookInline{b= } \textit{ac}, then \BookInline{a} \textbf{divides} \BookInline{b(}or \BookInline{a} \textbf{is a factor of} \BookInline{b),\allowbreak } denoted by \BookInline{a|b.} We read \BookInline{a\nmid b} as “\BookInline{a} does not divide \BookInline{b.}” \BookEnd{\blacksquare}

\BookLabel{34.2 Definition} An element \BookInline{u} of a commutative ring with unity \BookInline{R} is a\textbf{ unit of} \BookInline{R} if \BookInline{u} divides 1, that is, if \BookInline{u} has a multiplicative inverse in \BookInline{R.} Two elements \BookInline{a,\allowbreak b\in R} are\textbf{ associates in} \BookInline{R} if \BookInline{a= bu,\allowbreak } where \BookInline{u} is a unit in \BookInline{R.}

Exercise 27 asks us to show that this criterion for \BookInline{a} and \BookInline{b} to be associates is an equivalence relation on \BookInline{R.}\BookEnd{\blacksquare}

\BookLabel{34.3 Example} The only units in \BookInline{\mathbb{Z}} are 1 and \BookInline{- 1.} Thus the only associates of 26 in \BookInline{\mathbb{Z}} are 26 and \BookInline{- 26.}\BookEnd{\blacktriangle}

\BookLabel{34.4 Definition} A nonzero element \BookInline{p} that is not a unit of an integral domain \BookInline{D} is an\textbf{ irreducible of} \BookInline{D} if every factorization \BookInline{p= ab} in \BookInline{D} has the property that either \BookInline{a} or \BookInline{b} is a unit\BookInline{.}\BookEnd{\blacksquare}

Note that an associate of an irreducible \BookInline{p} is again an irreducible, for if \BookInline{p= uc} for a unit \BookInline{u,\allowbreak } then any factorization of \BookInline{c} provides a factorization of \BookInline{p.}

\BookLabel{34.5 Definition} An integral domain \BookInline{D} is a\textbf{ unique factorization domain} (abbreviated UFD) if the following conditions are satisfied:

\BookExercise{1.}{Every element of \BookInline{D} that is neither 0 nor a unit can be factored into a product of a finite number of irreducibles.}

\BookExercise{2.}{If \BookInline{p_{1}\cdots p_{r}} and \BookInline{q_{1}\cdots q_{s}} are two factorizations of the same element of \BookInline{D} into irreducibles, then \BookInline{r= s} and the \BookInline{q_{j}} can be renumbered so that \BookInline{p_{i}} and \BookInline{q_{i}} are associates\BookInline{.}\BookEnd{\blacksquare}}

\BookLabel{34.6 Example} Theorem 28.21 shows that for a field \BookInline{F,\allowbreak F[x]} is a UFD. Also we know that \BookInline{\mathbb{Z}} is a UFD; we have made frequent use of this fact, although we have never proved it. For example, in \BookInline{\mathbb{Z}} we have

\begin{center}\begin{adjustbox}{max width=\linewidth}\BookInline{\displaystyle 24= (2)(2)(3)(2)= (- 2)(- 3)(2)(2).}\end{adjustbox}\end{center}

\noindent Here 2 and \BookInline{- 2} are associates, as are 3 and \BookInline{- 3.} Thus except for order and associates, the irreducible factors in these two factorizations of 24 are the same\BookInline{.}\BookEnd{\blacktriangle}

Recall that the\textit{ principal ideal} \BookInline{\langle a\rangle } of \BookInline{D} consists of all multiples of the element \BookInline{a.} After just one more definition we can describe what we wish to achieve in this section.

\BookLabel{34.7 Definition} An integral domain \BookInline{D} is a\textbf{ principal ideal domain} (abbreviated PID) if every ideal in \BookInline{D} is a principal ideal\BookInline{.}\BookEnd{\blacksquare}

We know that \BookInline{\mathbb{Z}} is a PID because every ideal is of the form \BookInline{n\mathbb{Z},\allowbreak } generated by some integer \BookInline{n.} Theorem 31.24 shows that if \BookInline{F} is a field, then \BookInline{F[x]} is a PID. Our purpose in this section is to prove two exceedingly important theorems:

\BookExercise{1.}{Every PID is a UFD. (Theorem 34.18)}

\BookExercise{2.}{If \BookInline{D} is a UFD, then \BookInline{D[x]} is a UFD. (Theorem 34.30)}

\end{BookBody}

% Source PDF page 294; original printed page 277
\BookSource{294}{277}
\BookNarrativetrue

\begin{BookBody}

\BookTextHeading{\ensuremath{\blacksquare} \textbf{HISTORICAL NOTE}  }

\noindent The question of unique factorization in an integral domain other than the integers was first raised in public in connection with the attempted proof by Gabriel Lame´ (1795–1870) of Fermat’s Last Theorem, the conjecture that \BookInline{x^{n}+ y^{n}= z^{n}} has no nontrivial integral solutions for \BookInline{n>2.} It is not hard to show that the conjecture is true if it can be proved for all odd primes \BookInline{p.} At a meeting of the Paris Academy on March 1, 1847, Lame´ announced that he had proved the theorem and presented a sketch of the proof. Lame´’s idea was first to factor \BookInline{x^{p}+ y^{p}} over the complex numbers as

\begin{center}\begin{adjustbox}{max width=\linewidth}\BookInline{\displaystyle x^{p}+ y^{p}=}\end{adjustbox}\end{center}

\begin{center}\begin{adjustbox}{max width=\linewidth}\BookInline{\displaystyle (x+ y)(x+ \alpha y)(x+ \alpha ^{2}y)\cdots (x+ \alpha ^{p- 1}y)}\end{adjustbox}\end{center}

\noindent where \BookInline{\alpha } is a primitive \BookInline{p}th root of unity. He next proposed to show that if the factors in this expression are relatively prime and if \BookInline{x^{p}+ y^{p}= z^{p},\allowbreak } then each of the \BookInline{p} factors must be a \BookInline{p}th power. He could then demonstrate that this Fermat equation would be true for a triple \BookInline{x^{\prime },\allowbreak y^{\prime },\allowbreak z^{\prime },\allowbreak } each number smaller than the corresponding number in the original triple. This would lead to an infinite descending sequence of positive integers, an impossibility that would prove the theorem.

\noindent After Lame´ finished his announcement, however, Joseph Liouville (1809–1882) cast serious doubts on the purported proof, noting that the conclusion that each of the relatively prime factors was a \BookInline{p}th power because their product was a \BookInline{p}th power depended on the result that any integer can be uniquely factored into a product of primes. It was by no means clear that “integers” of the form \BookInline{x+ \alpha ^{k}y} had this unique factorization property. Although Lame´ attempted to overcome Liouville’s objections, the matter was settled on May 24, when Liouville produced a letter from Ernst Kummer noting that in 1844 he had already proved that unique factorization failed in the domain \BookInline{\mathbb{Z}[\alpha ],\allowbreak } where \BookInline{\alpha } is a 23rd root of unity.

It was not until 1994 that Fermat’s Last Theorem was proved, and by techniques of algebraic geometry unknown to Lame´ and Kummer. In the late 1950s, Yutaka Taniyama and Goro Shimura noticed a curious relationship between two seemingly disparate fields of mathematics, elliptic curves and modular forms. A few years after Taniyama’s tragic death at age 31, Shimura clarified this idea and eventually formulated what became known as the Taniyama–Shimura Conjecture. In 1984, Gerhard Frey asserted and in 1986 Ken Ribet proved that the Taniyama–Shimura Conjecture would imply the truth of Fermat’s Last Theorem. But it was finally Andrew Wiles of Princeton University who, after secretly working on this problem for seven years, gave a series of lectures at Cambridge University in June 1993 in which he announced a proof of enough of the Taniyama–Shimura Conjecture to derive Fermat’s Last Theorem. Unfortunately, a gap in the proof was soon discovered, and Wiles went back to work. It took him more than a year, but with the assistance of his student Richard Taylor, he finally was able to fill the gap. The result was published in the\textit{ Annals of Mathematics} in May 1995, and this 350-year-old problem was now solved.

The fact that \BookInline{F[x]} is a UFD, where \BookInline{F} is a field (by Theorem 28.21), illustrates both theorems. For by Theorem \BookInline{31.24,\allowbreak F[x]} is a PID. Also, since \BookInline{F} has no nonzero elements that are not units\BookInline{,\allowbreak F} satisfies our definition for a UFD. Thus Theorem 34.30 would give another proof that \BookInline{F[x]} is a UFD, except for the fact that we shall actually use Theorem 28.21 in proving Theorem 34.30. In the following section we shall study properties of a certain special class of UFDs, the\textit{ Euclidean domains}.

Let us proceed to prove the two theorems.

\BookTextHeading{\textbf{Every PID Is a UFD}  }

\noindent The steps leading up to Theorem 28.21 and its proof indicate the way for our proof of Theorem 34.18. Much of the material will be repetitive. We inefficiently handled the special case of \BookInline{F[x]} separately in Theorem 28.21, since it was easy and was the only case we needed for our field theory in general.

\end{BookBody}

% Source PDF page 295; original printed page 278
\BookSource{295}{278}
\BookNarrativetrue

\begin{BookBody}

To prove that an integral domain \BookInline{D} is a UFD, it is necessary to show that both Conditions 1 and 2 of the definition of a UFD are satisfied. For our special case of \BookInline{F[x]} in Theorem 28.21, Condition 1 was very easy and resulted from an argument that in a factorization of a polynomial of degree \BookInline{>0} into a product of two nonconstant polynomials, the degree of each factor was less than the degree of the original polynomial. Thus we couldn’t keep on factoring indefinitely without running into unit factors, that is, polynomials of degree 0. For the general case of a PID, it is harder to show that this is so. We now turn to this problem. We shall need the definition of the union of an arbitrary collection of sets. The definition must include the possibility that the collection of sets is infinite.

\BookLabel{34.8 Definition} If \BookInline{\{A_{i}\mid i\in I\}} is a collection of sets, then the\textbf{ union} \BookInline{\cup _{i\in I}A_{i}} \textbf{of the sets} \BookInline{A_{i}} is the set of all \BookInline{x} such that \BookInline{x\in A_{i}} for at least one \BookInline{i\in I.}\BookEnd{\blacksquare}

\BookLabel{34.9 Lemma} Let \BookInline{R} be a commutative ring and let \BookInline{N_{1}\subseteq N_{2}\subseteq \cdots } be an ascending chain of ideals \BookInline{N_{i}} in \BookInline{R.} Then \BookInline{N= \cup _{i}N_{i}} is an ideal of \BookInline{R.}

\BookLabel{\textit{Proof}} Let \BookInline{a,\allowbreak b\in N.} Then there are ideals \BookInline{N_{i}} and \BookInline{N_{j}} in the chain, with \BookInline{a\in N_{i}} and \BookInline{b\in N_{j}.} Now either \BookInline{N_{i}\subseteq N_{j}} or \BookInline{N_{j}\subseteq N_{i};} let us assume that \BookInline{N_{i}\subseteq N_{j},\allowbreak } so both \BookInline{a} and \BookInline{b} are in \BookInline{N_{j}.} This implies that \BookInline{a\pm b} and \BookInline{ab} are in \BookInline{N_{j},\allowbreak } so \BookInline{a\pm b} and \BookInline{ab} are in \BookInline{N.} Taking \BookInline{a= 0,\allowbreak } we see that \BookInline{b\in N} implies \BookInline{- b\in N,\allowbreak } and \BookInline{0\in N} since \BookInline{0\in N_{i}.} Thus \BookInline{N} is a subring of \BookInline{D.} For \BookInline{a\in N} and \BookInline{d\in D,\allowbreak } we must have \BookInline{a\in N_{i}} for some \BookInline{N_{i}.} Then since \BookInline{N_{i}} is an ideal\BookInline{,\allowbreak da= ad} is in \BookInline{N_{i}.} Therefore\BookInline{,\allowbreak da\in \cup _{i}N_{i},\allowbreak } that is\BookInline{,\allowbreak da\in N.} Hence \BookInline{N} is an ideal\BookInline{.}\BookEnd{\blacklozenge}

\BookLabel{34.10 Lemma} \textbf{(Ascending Chain Condition for a PID)} Let \BookInline{D} be a PID. If \BookInline{N_{1}\subseteq N_{2}\subseteq \cdots } is an ascending chain of ideals \BookInline{N_{i},\allowbreak } then there exists a positive integer \BookInline{r} such that \BookInline{N_{r}= N_{s}} for all \BookInline{s\geq r.} Equivalently, every strictly ascending chain of ideals (all inclusions proper) in a PID is of finite length. We express this by saying that the\textbf{ ascending chain condition} (ACC) holds for ideals in a PID.

\BookLabel{\textit{Proof}} By Lemma 34.9, we know that \BookInline{N= \cup _{i}N_{i}} is an ideal of \BookInline{D.} Now as an ideal in \BookInline{D,\allowbreak } which is a PID\BookInline{,\allowbreak N= \langle c\rangle } for some \BookInline{c\in D.} Since \BookInline{N= \cup _{i}N_{i},\allowbreak } we must have \BookInline{c\in N_{r},\allowbreak } for some \BookInline{r\in \mathbb{Z}^{+ }.} For \BookInline{s\geq r,\allowbreak } we have

\begin{center}\begin{adjustbox}{max width=\linewidth}\BookInline{\displaystyle \langle c\rangle \subseteq N_{r}\subseteq N_{s}\subseteq N= \langle c\rangle .}\end{adjustbox}\end{center}

\noindent Thus \BookInline{N_{r}= N_{s}} for \BookInline{s\geq r.}

The equivalence with the ACC is immediate. \BookEnd{\blacklozenge}

\BookLabel{34.11 Definition} A commutative ring with unity \BookInline{R} that satisfies the ascending chain condition is a \textbf{Noetherian ring}. That is, a commutative ring with unity \BookInline{R} is Noetherian if for every chain of ideals \BookInline{N_{1}\subseteq N_{2}\subseteq N_{3}\subseteq...} in \BookInline{R,\allowbreak } there is an integer \BookInline{r} such that if \BookInline{s\geq r,\allowbreak } then \BookInline{N_{r}= N_{s}.}\BookEnd{\blacksquare}

\noindent Lemma 34.10 states that every PID is a Noetherian ring. In Section 37 we will see that if \BookInline{R} is a Noetherian ring, then \BookInline{R[x]} is also a Noetherian ring.

In what follows, it will be useful to remember that for elements \BookInline{a} and \BookInline{b} of a domain \BookInline{D,\allowbreak }

\BookInline{\langle a\rangle \subseteq \langle b\rangle } if and only if \BookInline{b} divides \BookInline{a,\allowbreak } and

\BookInline{\langle a\rangle = \langle b\rangle } if and only if \BookInline{a} and \BookInline{b} are associates.

\noindent For the first property, note that \BookInline{\langle a\rangle \subseteq \langle b\rangle } if and only if \BookInline{a\in \langle b\rangle,\allowbreak } which is true if and only if \BookInline{a= bd} for some \BookInline{d\in D,\allowbreak } so that \BookInline{b} divides \BookInline{a.} Using this first property, we see that \BookInline{\langle a\rangle = \langle b\rangle } if and only if \BookInline{a= bc} and \BookInline{b= ad} for some \BookInline{c,\allowbreak d\in D.} But then \BookInline{a= adc} and by canceling, we obtain \BookInline{1= } \textit{dc}. Thus \BookInline{d} and \BookInline{c} are units, so \BookInline{a} and \BookInline{b} are associates.

\end{BookBody}

% Source PDF page 296; original printed page 279
\BookSource{296}{279}
\BookNarrativetrue

\begin{BookBody}

We can now prove Condition 1 of the definition of a UFD for an integral domain that is a PID.

\BookLabel{34.12 Theorem} Let \BookInline{D} be a PID. Every element that is neither 0 nor a unit in \BookInline{D} is a product of irreducibles.

\BookLabel{\textit{Proof}} Let \BookInline{a\in D,\allowbreak } where \BookInline{a} is neither 0 nor a unit. We first show that \BookInline{a} has at least one irreducible factor. If \BookInline{a} is an irreducible, we are done. If \BookInline{a} is not an irreducible, then \BookInline{a= a_{1}b_{1},\allowbreak } where \textit{neither} \BookInline{a_{1}} nor \BookInline{b_{1}} is a unit. Now

\begin{center}\begin{adjustbox}{max width=\linewidth}\BookInline{\displaystyle \langle a\rangle \subset \langle a_{1}\rangle ,\allowbreak }\end{adjustbox}\end{center}

\noindent for \BookInline{\langle a\rangle \subseteq \langle a_{1}\rangle } follows from \BookInline{a= a_{1}b_{1},\allowbreak } and if \BookInline{\langle a\rangle = \langle a_{1}\rangle,\allowbreak } then \BookInline{a} and \BookInline{a_{1}} would be associates and \BookInline{b_{1}} would be a unit, contrary to construction. Continuing this procedure then, starting now with \BookInline{a_{1},\allowbreak } we arrive at a strictly ascending chain of ideals

\begin{center}\begin{adjustbox}{max width=\linewidth}\BookInline{\displaystyle \langle a\rangle \subset \langle a_{1}\rangle \subset \langle a_{2}\rangle \subset \cdots .}\end{adjustbox}\end{center}

\noindent By the ACC in Lemma 34.10, this chain terminates with some \BookInline{\langle a_{r}\rangle,\allowbreak } and \BookInline{a_{r}} must then be irreducible. Thus \BookInline{a} has an irreducible factor \BookInline{a_{r}.}

By what we have just proved, for an element \BookInline{a} that is neither 0 nor a unit in \BookInline{D,\allowbreak } either \BookInline{a} is irreducible or \BookInline{a= p_{1}c_{1}} for \BookInline{p_{1}} an irreducible and \BookInline{c_{1}} not a unit. By an argument similar to the one just made, in the latter case we can conclude that \BookInline{\langle a\rangle \subset \langle c_{1}\rangle.} If \BookInline{c_{1}} is not irreducible, then \BookInline{c_{1}= p_{2}c_{2}} for an irreducible \BookInline{p_{2}} with \BookInline{c_{2}} not a unit. Continuing, we get a strictly ascending chain of ideals

\begin{center}\begin{adjustbox}{max width=\linewidth}\BookInline{\displaystyle \langle a\rangle \subset \langle c_{1}\rangle \subset \langle c_{2}\rangle \subset \cdots .}\end{adjustbox}\end{center}

\noindent This chain must terminate, by the ACC in Lemma 34.10, with some \BookInline{c_{r}= q_{r}} that is an irreducible. Then \BookInline{a= p_{1}p_{2}\cdots p_{r}q_{r}.}\BookEnd{\blacklozenge}

This completes our demonstration of Condition 1 of the definition of a UFD. Let us turn to Condition 2. Our arguments here are parallel to those leading to Theorem 28.21. The results we encounter along the way are of some interest in themselves.

\BookLabel{34.13 Lemma} \textbf{(Generalization of Theorem 31.25)} An ideal \BookInline{\langle p\rangle } in a PID is maximal if and only if \BookInline{p} is an irreducible.

\BookLabel{\textit{Proof}} Let \BookInline{\langle p\rangle } be a maximal ideal of \BookInline{D,\allowbreak } a PID. Suppose that \BookInline{p= ab} in \BookInline{D.} Then \BookInline{\langle p\rangle \subseteq \langle a\rangle.} Suppose that \BookInline{\langle a\rangle = \langle p\rangle.} Then \BookInline{a} and \BookInline{p} would be associates, so \BookInline{b} must be a unit. If \BookInline{\langle a\rangle \neq } \BookInline{\langle p\rangle,\allowbreak } then we must have \BookInline{\langle a\rangle = \langle 1\rangle = D,\allowbreak } since \BookInline{\langle p\rangle } is maximal. But then \BookInline{a} and 1 are associates, so \BookInline{a} is a unit. Thus, if \BookInline{p= ab,\allowbreak } either \BookInline{a} or \BookInline{b} must be a unit. Hence \BookInline{p} is an irreducible of \BookInline{D.}

Conversely, suppose that \BookInline{p} is an irreducible in \BookInline{D.} Then if \BookInline{\langle p\rangle \subseteq \langle a\rangle,\allowbreak } we must have \BookInline{p= ab.} Now if \BookInline{a} is a unit, then \BookInline{\langle a\rangle = \langle 1\rangle = D.} If \BookInline{a} is not a unit, then \BookInline{b} must be a unit, so there exists \BookInline{u\in D} such that \BookInline{bu= 1.} Then \BookInline{pu= } \textit{abu} \BookInline{= a,\allowbreak } so \BookInline{\langle a\rangle \subseteq \langle p\rangle,\allowbreak } and we have \BookInline{\langle a\rangle = \langle p\rangle.} Thus \BookInline{\langle p\rangle \subseteq \langle a\rangle } implies that either \BookInline{\langle a\rangle = D} or \BookInline{\langle a\rangle = \langle p\rangle,\allowbreak } and \BookInline{\langle p\rangle \neq D} or \BookInline{p} would be a unit. Hence \BookInline{\langle p\rangle } is a maximal ideal\BookInline{.}\BookEnd{\blacklozenge}

\BookLabel{34.14 Lemma} \textbf{(Generalization of Theorem 31.27)} In a PID, if an irreducible \BookInline{p} divides \BookInline{ab,\allowbreak } then either \BookInline{p|a} or \BookInline{p|b.}

\BookLabel{\textit{Proof}} Let \BookInline{D} be a PID and suppose that for an irreducible \BookInline{p} in \BookInline{D} we have \BookInline{p|ab.} Then \BookInline{(ab)\in } \BookInline{\langle p\rangle.} Since every maximal ideal in \BookInline{D} is a prime ideal by Corollary \BookInline{31.16,\allowbreak (ab)\in \langle p\rangle } implies that either \BookInline{a\in \langle p\rangle } or \BookInline{b\in \langle p\rangle,\allowbreak } giving either \BookInline{p|a} or \BookInline{p|b.}\BookEnd{\blacklozenge}

\BookLabel{34.15 Corollary} If \BookInline{p} is an irreducible in a PID and \BookInline{p} divides the product \BookInline{a_{1}a_{2}\cdots a_{n}} for \BookInline{a_{i}\in D,\allowbreak } then \BookInline{p|a_{i}} for at least one \BookInline{i.}

\end{BookBody}

% Source PDF page 297; original printed page 280
\BookSource{297}{280}
\BookNarrativetrue

\begin{BookBody}

\BookLabel{\textit{Proof}} Proof of this corollary is immediate from Lemma 34.14 if we use mathematical induction\BookInline{.}\BookEnd{\blacklozenge}

\BookLabel{34.16 Definition} A nonzero nonunit element \BookInline{p} of an integral domain \BookInline{D} is a\textbf{ prime} if, for all \BookInline{a,\allowbreak b\in D,\allowbreak } \BookInline{p|ab} implies either \BookInline{p|a} or \BookInline{p|b.}\BookEnd{\blacksquare}

Lemma 34.14 focused our attention on the defining property of a prime. In Exercises 25 and 26, we ask you to show that a prime in an integral domain is always an irreducible and that in a UFD an irreducible is also a prime. Thus the concepts of prime and irreducible coincide in a UFD. Example 34.17 will exhibit an integral domain containing some irreducibles that are not primes, so the concepts do not coincide in every domain.

\BookLabel{34.17 Example} Let \BookInline{F} be a field and let \BookInline{D} be the subdomain \BookInline{F[x^{3},\allowbreak xy,\allowbreak y^{3}]} of \BookInline{F[x,\allowbreak y].} Then \BookInline{x^{3},\allowbreak xy,\allowbreak } and \BookInline{y^{3}} are irreducibles in \BookInline{D,\allowbreak } but

\begin{center}\begin{adjustbox}{max width=\linewidth}\BookInline{\displaystyle (x^{3})(y^{3})= (xy)(xy)(xy).}\end{adjustbox}\end{center}

\noindent Since \BookInline{xy} divides \BookInline{x^{3}y^{3}} but not \BookInline{x^{3}} or \BookInline{y^{3},\allowbreak } we see that \BookInline{xy} is not a prime. Similar arguments show that neither \BookInline{x^{3}} nor \BookInline{y^{3}} is a prime\BookInline{.}\BookEnd{\blacktriangle}

The defining property of a prime is precisely what is needed to establish uniqueness of factorization, Condition 2 in the definition of a UFD. We now complete the proof of Theorem 34.18 by demonstrating the uniqueness of factorization in a PID.

\BookTextHeading{\textbf{(Generalization of Theorem 28.21)} Every PID is a UFD.   }

\BookLabel{34.18 Theorem} 

\BookLabel{\textit{Proof}} Theorem 34.12 shows that if \BookInline{D} is a PID, then each \BookInline{a\in D,\allowbreak } where \BookInline{a} is neither 0 nor a unit, has a factorization

\begin{center}\begin{adjustbox}{max width=\linewidth}\BookInline{\displaystyle a= p_{1}p_{2}\cdots p_{r}}\end{adjustbox}\end{center}

\noindent into irreducibles. It remains for us to show uniqueness. Let

\begin{center}\begin{adjustbox}{max width=\linewidth}\BookInline{\displaystyle a= q_{1}q_{2}\cdots q_{s}}\end{adjustbox}\end{center}

\noindent be another such factorization into irreducibles. Then we have \BookInline{p_{1}|(q_{1}q_{2}\cdots q_{s}),\allowbreak } which implies that \BookInline{p_{1}|q_{j}} for some \BookInline{j} by Corollary 34.15. By changing the order of the \BookInline{q_{j}} if necessary, we can assume that \BookInline{j= 1} so \BookInline{p_{1}|q_{1}.} Then \BookInline{q_{1}= p_{1}u_{1},\allowbreak } and since \BookInline{q_{1}} is an irreducible\BookInline{,\allowbreak u_{1}} is a unit, so \BookInline{p_{1}} and \BookInline{q_{1}} are associates. We have then

\begin{center}\begin{adjustbox}{max width=\linewidth}\BookInline{\displaystyle p_{1}p_{2}\cdots p_{r}= p_{1}u_{1}q_{2}\cdots q_{s},\allowbreak }\end{adjustbox}\end{center}

\noindent so by the cancellation law in \BookInline{D,\allowbreak }

\begin{center}\begin{adjustbox}{max width=\linewidth}\BookInline{\displaystyle p_{2}\cdots p_{r}= u_{1}q_{2}\cdots q_{s}.}\end{adjustbox}\end{center}

\noindent Continuing this process, starting with \BookInline{p_{2}} and so on, we finally arrive at

\begin{center}\begin{adjustbox}{max width=\linewidth}\BookInline{\displaystyle 1= u_{1}u_{2}\cdots u_{r}q_{r+ 1}\cdots q_{s}.}\end{adjustbox}\end{center}

\noindent Since the \BookInline{q_{j}} are irreducibles, we must have \BookInline{r= s.} \BookEnd{\blacklozenge}

Example 34.32 at the end of this section will show that the converse to Theorem 34.18 is false. That is, a UFD need not be a PID.

Many algebra texts start by proving the following corollary of Theorem 34.18. We have assumed that you were familiar with this corollary and used it freely in our other work.

\BookLabel{34.19 Corollary} \textbf{(Fundamental Theorem of Arithmetic)} The integral domain \BookInline{\mathbb{Z}} is a UFD.

\end{BookBody}

% Source PDF page 298; original printed page 281
\BookSource{298}{281}
\BookNarrativetrue

\begin{BookBody}

\BookLabel{\textit{Proof}} We have seen that all ideals in \BookInline{\mathbb{Z}} are of the form \BookInline{n\mathbb{Z}= \langle n\rangle } for \BookInline{n\in \mathbb{Z}.} Thus \BookInline{\mathbb{Z}} is a PID, and Theorem 34.18 applies\BookInline{.}\BookEnd{\blacklozenge}

It is worth noting that the proof that \BookInline{\mathbb{Z}} is a PID was really way back in Corollary 6.7. We proved Theorem 6.6 by using the division algorithm for \BookInline{\mathbb{Z}} exactly as we proved, in Theorem 31.24, that \BookInline{F[x]} is a PID by using the division algorithm for \BookInline{F[x].} In Section 35, we shall examine this parallel more closely.

\BookTextHeading{\textbf{If} \textit{\textbf{D}} \textbf{Is a UFD, Then} \textit{\textbf{D}}\textbf{[}\textit{\textbf{x}}\textbf{] Is a UFD}     }

\noindent We now start the proof of Theorem 34.30, our second main result for this section. The idea of the argument is as follows. Let \BookInline{D} be a UFD. We can form a field of quotients \BookInline{F} of \BookInline{D.} Then \BookInline{F[x]} is a UFD by Theorem 28.21, and we shall show that we can recover a factorization for \BookInline{f(x)\in D[x]} from its factorization in \BookInline{F[x].} It will be necessary to compare the irreducibles in \BookInline{F[x]} with those in \BookInline{D[x],\allowbreak } of course. This approach, which we prefer as more intuitive than some more efficient modern ones, is essentially due to Gauss.

\BookLabel{34.20 Definition} Let \BookInline{D} be a UFD and let \BookInline{a_{1},\allowbreak a_{2},\allowbreak \cdots,\allowbreak a_{n}} be nonzero elements of \BookInline{D.} An element \BookInline{d} of \BookInline{D} is a\textbf{ greatest common divisor} (abbreviated gcd)\textbf{ of all of the} \BookInline{a_{i}} if \BookInline{d|a_{i}} for \BookInline{i= 1,\allowbreak \cdots,\allowbreak n} and any other \BookInline{d^{\prime }\in D} that divides all the \BookInline{a_{i}} also divides \BookInline{d.}\BookEnd{\blacksquare}

In this definition, we called \BookInline{d} “a” gcd rather than “the” gcd because gcd’s are only defined up to units. Suppose that \BookInline{d} and \BookInline{d^{\prime }} are two gcd’s of \BookInline{a_{i}} for \BookInline{i= 1,\allowbreak \cdots,\allowbreak n.} Then \BookInline{d|d^{\prime }} and \BookInline{d^{\prime }|d} by our definition. Thus \BookInline{d= q^{\prime }d^{\prime }} and \BookInline{d^{\prime }= qd} for some \BookInline{q,\allowbreak q^{\prime }\in D,\allowbreak } so \BookInline{1d= q^{\prime }qd.} By cancellation in \BookInline{D,\allowbreak } we see that \BookInline{q^{\prime }q= 1} so \BookInline{q} and \BookInline{q^{\prime }} are indeed units.

The technique in the example that follows shows that gcd’s exist in a UFD.

\BookLabel{34.21 Example} Let us find a gcd of \BookInline{420,\allowbreak - 168,\allowbreak } and 252 in the UFD \BookInline{\mathbb{Z}.} Factoring, we obtain \BookInline{420= 2^{2}\cdot } \BookInline{3\cdot 5\cdot 7,\allowbreak - 168= 2^{3}\cdot (- 3)\cdot 7,\allowbreak } and \BookInline{252= 2^{2}\cdot 3^{2}\cdot 7.} We choose one of these numbers, say 420, and find the highest power of each of its irreducible factors (up to associates) that divides all the numbers\BookInline{,\allowbreak 420,\allowbreak - 168,\allowbreak } and 252 in our case. We take as gcd the product of these highest powers of irreducibles. For our example, these powers of irreducible factors of 420 are 22, 31, 50, and 71 so we take as gcd \BookInline{d= 4\cdot 3\cdot 1\cdot 7= 84.} The only other gcd of these numbers in \BookInline{\mathbb{Z}} is \BookInline{- 84,\allowbreak } because 1 and \BookInline{- 1} are the only units\BookInline{.}\BookEnd{\blacktriangle}

Execution of the technique in Example 34.21 depends on being able to factor an element of a UFD into a product of irreducibles. This can be a tough job, even in \BookInline{\mathbb{Z}.} Section 35 will exhibit a technique, the Euclidean Algorithm, that will allow us to find gcd’s without factoring in a class of UFD’s that includes \BookInline{\mathbb{Z}} and \BookInline{F[x]} for a field \BookInline{F.}

\BookLabel{34.22 Definition} Let \BookInline{D} be a UFD. A nonconstant polynomial

\begin{center}\begin{adjustbox}{max width=\linewidth}\BookInline{\displaystyle f(x)= a_{0}+ a_{1}x+ \cdots + a_{n}x^{n}}\end{adjustbox}\end{center}

\noindent in \BookInline{D[x]} is\textbf{ primitive} if 1 is a gcd of the \BookInline{a_{i}} for \BookInline{i= 0,\allowbreak 1,\allowbreak \cdots,\allowbreak n.} \BookEnd{\blacksquare}

\BookLabel{34.23 Example} In \BookInline{\mathbb{Z}[x],\allowbreak 4x^{2}+ 3x+ 2} is primitive, but \BookInline{4x^{2}+ 6x+ 2} is not, since 2, a nonunit in \BookInline{\mathbb{Z},\allowbreak } is a common divisor of 4, 6, and \BookInline{2.}\BookEnd{\blacktriangle}

Observe that every nonconstant irreducible in \BookInline{D[x]} must be a primitive polynomial.

\BookLabel{34.24 Lemma} If \BookInline{D} is a UFD, then for every nonconstant \BookInline{f(x)\in D[x]} we have \BookInline{f(x)= (c)g(x),\allowbreak } where \BookInline{c\in D,\allowbreak g(x)\in D[x],\allowbreak } and \BookInline{g(x)} is primitive. The element \BookInline{c} is unique up to a unit factor in \BookInline{D} and is the\textbf{ content of} \BookInline{f(x).} Also \BookInline{g(x)} is unique up to a unit factor in \BookInline{D.}

\end{BookBody}

% Source PDF page 299; original printed page 282
\BookSource{299}{282}
\BookNarrativetrue

\begin{BookBody}

\BookLabel{\textit{Proof}} Let \BookInline{f(x)\in D[x]} be given where \BookInline{f(x)} is a nonconstant polynomial with coefficients \BookInline{a_{0},\allowbreak a_{1},\allowbreak \cdots,\allowbreak a_{n}.} Let \BookInline{c} be a gcd of the \BookInline{a_{i}} for \BookInline{i= 0,\allowbreak 1,\allowbreak \cdots,\allowbreak n.} Then for each \BookInline{i,\allowbreak } we have \BookInline{a_{i}= cq_{i}} for some \BookInline{q_{i}\in D.} By the distributive law, we have \BookInline{f(x)= (c)g(x),\allowbreak } where no irreducible in \BookInline{D} divides all of the coefficients \BookInline{q_{0},\allowbreak q_{1},\allowbreak \cdots,\allowbreak q_{n}} of \BookInline{g(x).} Thus \BookInline{g(x)} is a primitive polynomial.

For uniqueness, if also \BookInline{f(x)= (d)h(x)} for \BookInline{d\in D,\allowbreak h(x)\in D[x],\allowbreak } and \BookInline{h(x)} primitive, then each irreducible factor of \BookInline{c} must divide \BookInline{d} and conversely. By setting \BookInline{(c)g(x)= } \BookInline{(d)h(x)} and canceling irreducible factors of \BookInline{c} into \BookInline{d,\allowbreak } we arrive at \BookInline{(u)g(x)= (v)h(x)} for a unit \BookInline{u\in D.} But then \BookInline{v} must be a unit of \BookInline{D} or we would be able to cancel irreducible factors of \BookInline{v} into \BookInline{u.} Thus \BookInline{u} and \BookInline{v} are both units, so \BookInline{c} is unique up to a unit factor. From \BookInline{f(x)= } \BookInline{(c)g(x),\allowbreak } we see that the primitive polynomial \BookInline{g(x)} is also unique up to a unit factor\BookInline{.}\BookEnd{\blacklozenge}

\BookLabel{34.25 Example} In \BookInline{\mathbb{Z}[x],\allowbreak }

\begin{center}\begin{adjustbox}{max width=\linewidth}\BookInline{\displaystyle 4x^{2}+ 6x- 8= (2)(2x^{2}+ 3x- 4),\allowbreak }\end{adjustbox}\end{center}

\noindent where \BookInline{2x^{3}+ 3x- 4} is primitive. \BookEnd{\blacktriangle}

\BookLabel{34.26 Lemma} \textbf{(Gauss’s Lemma)} If \BookInline{D} is a UFD, then a product of two primitive polynomials in \BookInline{D[x]} is again primitive.

\BookLabel{\textit{Proof}} Let

\begin{center}\begin{adjustbox}{max width=\linewidth}\BookInline{\displaystyle f(x)= a_{0}+ a_{1}x+ \cdots + a_{n}x^{n}}\end{adjustbox}\end{center}

\noindent and

\begin{center}\begin{adjustbox}{max width=\linewidth}\BookInline{\displaystyle g(x)= b_{0}+ b_{1}x+ \cdots + b_{m}x^{m}}\end{adjustbox}\end{center}

\noindent be primitive in \BookInline{D[x],\allowbreak } and let \BookInline{h(x)= f(x)g(x).} Let \BookInline{p} be an irreducible in \BookInline{D.} Then \BookInline{p} does not divide all \BookInline{a_{i}} and \BookInline{p} does not divide all \BookInline{b_{j},\allowbreak } since \BookInline{f(x)} and \BookInline{g(x)} are primitive. Let \BookInline{a_{r}} be the first coefficient of \BookInline{f(x)} not divisible by \BookInline{p;} that is\BookInline{,\allowbreak p|a_{i}} for \BookInline{i<r,\allowbreak } but \BookInline{p\nmid a_{r}(}that is, \BookInline{p} does not divide \BookInline{a_{r}).} Similarly, let \BookInline{p|b_{j}} for \BookInline{j<s,\allowbreak } but \BookInline{p\nmid b_{s}.} The coefficient of \BookInline{x^{r+ s}} in \BookInline{h(x)= f(x)g(x)} is

\begin{center}\begin{adjustbox}{max width=\linewidth}\BookInline{\displaystyle c_{r+ s}= (a_{0}b_{r+ s}+ \cdots + a_{r- 1}b_{s+ 1})+ a_{r}b_{s}+ (a_{r+ 1}b_{s- 1}+ \cdots + a_{r+ s}b_{0}).}\end{adjustbox}\end{center}

\noindent Now \BookInline{p|a_{i}} for \BookInline{i<r} implies that

\begin{center}\begin{adjustbox}{max width=\linewidth}\BookInline{\displaystyle p|(a_{0}b_{r+ s}+ \cdots + a_{r- 1}b_{s+ 1}),\allowbreak }\end{adjustbox}\end{center}

\noindent and also \BookInline{p|b_{j}} for \BookInline{j<s} implies that

\begin{center}\begin{adjustbox}{max width=\linewidth}\BookInline{\displaystyle p|(a_{r+ 1}b_{s- 1}+ \cdots + a_{r+ s}b_{0}).}\end{adjustbox}\end{center}

\noindent But \BookInline{p} does not divide \BookInline{a_{r}} or \BookInline{b_{s},\allowbreak } so \BookInline{p} does not divide \BookInline{a_{r}b_{s},\allowbreak } and consequently \BookInline{p} does not divide \BookInline{c_{r+ s}.} This shows that given an irreducible \BookInline{p\in D,\allowbreak } there is some coefficient of \BookInline{f(x)g(x)} not divisible by \BookInline{p.} Thus \BookInline{f(x)g(x)} is primitive\BookInline{.}\BookEnd{\blacklozenge}

\BookLabel{34.27 Corollary} If \BookInline{D} is a UFD, then a finite product of primitive polynomials in \BookInline{D[x]} is again primitive.

\BookLabel{\textit{Proof}} This corollary follows from Lemma 34.26 by induction. \BookEnd{\blacklozenge}

Now let \BookInline{D} be a UFD and let \BookInline{F} be a field of quotients of \BookInline{D.} By Theorem 28.21, \BookInline{F[x]} is a UFD. As we said earlier, we shall show that \BookInline{D[x]} is a UFD by carrying a factorization in \BookInline{F[x]} of \BookInline{f(x)\in D[x]} back into one in \BookInline{D[x].} The next lemma relates the nonconstant irreducibles of \BookInline{D[x]} to those of \BookInline{F[x].} This is the last important step.

\BookLabel{34.28 Lemma} Let \BookInline{D} be a UFD and let \BookInline{F} be a field of quotients of \BookInline{D.} Let \BookInline{f(x)\in D[x],\allowbreak } where (degree \BookInline{f(x))>0.} If \BookInline{f(x)} is an irreducible in \BookInline{D[x],\allowbreak } then \BookInline{f(x)} is also an irreducible in \BookInline{F[x].} Also, if \BookInline{f(x)} is primitive in \BookInline{D[x]} and irreducible in \BookInline{F[x],\allowbreak } then \BookInline{f(x)} is irreducible in \BookInline{D[x].}

\end{BookBody}

% Source PDF page 300; original printed page 283
\BookSource{300}{283}
\BookNarrativetrue

\begin{BookBody}

\BookLabel{\textit{Proof}} Suppose that a nonconstant \BookInline{f(x)\in D[x]} factors into polynomials of lower degree in \BookInline{F[x],\allowbreak } that is,

\begin{center}\begin{adjustbox}{max width=\linewidth}\BookInline{\displaystyle f(x)= r(x)s(x)}\end{adjustbox}\end{center}

\noindent for \BookInline{r(x),\allowbreak s(x)\in F[x].} Then since \BookInline{F} is a field of quotients of \BookInline{D,\allowbreak } each coefficient in \BookInline{r(x)} and \BookInline{s(x)} is of the form \BookInline{a/b} for some \BookInline{a,\allowbreak b\in D.} By clearing denominators, we can get

\begin{center}\begin{adjustbox}{max width=\linewidth}\BookInline{\displaystyle (d)f(x)= r_{1}(x)s_{1}(x)}\end{adjustbox}\end{center}

\noindent for \BookInline{d\in D,\allowbreak } and \BookInline{r_{1}(x),\allowbreak s_{1}(x)\in D[x],\allowbreak } where the degrees of \BookInline{r_{1}(x)} and \BookInline{s_{1}(x)} are the degrees of \BookInline{r(x)} and \BookInline{s(x),\allowbreak } respectively. By Lemma \BookInline{34.24,\allowbreak f(x)= (c)g(x),\allowbreak r_{1}(x)= (c_{1})r_{2}(x),\allowbreak } and \BookInline{s_{1}(x)= (c_{2})s_{2}(x)} for primitive polynomials \BookInline{g(x),\allowbreak r_{2}(x),\allowbreak } and \BookInline{s_{2}(x),\allowbreak } and \BookInline{c,\allowbreak c_{1},\allowbreak c_{2}\in D.} Then

\begin{center}\begin{adjustbox}{max width=\linewidth}\BookInline{\displaystyle (dc)g(x)= (c_{1}c_{2})r_{2}(x)s_{2}(x),\allowbreak }\end{adjustbox}\end{center}

\noindent and by Lemma \BookInline{34.26,\allowbreak r_{2}(x)s_{2}(x)} is primitive. By the uniqueness part of Lemma 34.24, \BookInline{c_{1}c_{2}= dcu} for some unit \BookInline{u} in \BookInline{D.} But then

\begin{center}\begin{adjustbox}{max width=\linewidth}\BookInline{\displaystyle (dc)g(x)= (dcu)r_{2}(x)s_{2}(x),\allowbreak }\end{adjustbox}\end{center}

\noindent so

\begin{center}\begin{adjustbox}{max width=\linewidth}\BookInline{\displaystyle f(x)= (c)g(x)= (cu)r_{2}(x)s_{2}(x).}\end{adjustbox}\end{center}

\noindent \textit{We have shown that if} \BookInline{f(x)} \textit{factors nontrivially in} \BookInline{F[x],\allowbreak } \textit{then} \BookInline{f(x)} \textit{factors nontrivially into} \textit{polynomials of the same degrees in} \BookInline{D[x].} Thus if \BookInline{f(x)\in D[x]} is irreducible in \BookInline{D[x],\allowbreak } it must be irreducible in \BookInline{F[x].}

A nonconstant \BookInline{f(x)\in D[x]} that is primitive in \BookInline{D[x]} and irreducible in \BookInline{F[x]} is also irreducible in \BookInline{D[x],\allowbreak } since \BookInline{D[x]\subseteq F[x].}\BookEnd{\blacklozenge}

Lemma 34.28 shows that if \BookInline{D} is a UFD, the irreducibles in \BookInline{D[x]} are precisely the irreducibles in \BookInline{D,\allowbreak } together with the nonconstant primitive polynomials that are irreducible in \BookInline{F[x],\allowbreak } where \BookInline{F} is a field of quotients of \BookInline{D[x].}

The preceding lemma is very important in its own right. This is indicated by the following corollary, a special case of which was our Theorem 28.12. (We admit that it does not seem very sensible to call a special case of a corollary of a lemma a theorem. The label assigned to a result depends somewhat on the context in which it appears.)

\BookLabel{34.29 Corollary} If \BookInline{D} is a UFD and \BookInline{F} is a field of quotients of \BookInline{D,\allowbreak } then a nonconstant \BookInline{f(x)\in D[x]} factors into a product of two polynomials of lower degrees \BookInline{r} and \BookInline{s} in \BookInline{F[x]} if and only if it has a factorization into polynomials of the same degrees \BookInline{r} and \BookInline{s} in \BookInline{D[x].}

\BookLabel{\textit{Proof}} It was shown in the proof of Lemma 34.28 that if \BookInline{f(x)} factors into a product of two polynomials of lower degree in \BookInline{F[x],\allowbreak } then it has a factorization into polynomials of the same degrees in \BookInline{D[x](}see the next-to-last sentence of the first paragraph of the proof).

The converse holds since \BookInline{D[x]\subseteq F[x].} \BookEnd{\blacklozenge}

\noindent We are now prepared to prove our main theorem.

\BookLabel{34.30 Theorem} If \BookInline{D} is a UFD, then \BookInline{D[x]} is a UFD.

\BookLabel{\textit{Proof}} Let \BookInline{f(x)\in D[x],\allowbreak } where \BookInline{f(x)} is neither 0 nor a unit. If \BookInline{f(x)} is of degree 0, we are done, since \BookInline{D} is a UFD. Suppose that (degree \BookInline{f(x))>0.} Let

\begin{center}\begin{adjustbox}{max width=\linewidth}\BookInline{\displaystyle f(x)= g_{1}(x)g_{2}(x)\cdots g_{r}(x)}\end{adjustbox}\end{center}

\noindent be a factorization of \BookInline{f(x)} in \BookInline{D[x]} having the greatest number \BookInline{r} of factors of positive degree. (There is such a greatest number of such factors because \BookInline{r} cannot exceed the degree of \BookInline{f(x).)} Now factor each \BookInline{g_{i}(x)} in the form \BookInline{g_{i}(x)= c_{i}h_{i}(x)} where \BookInline{c_{i}} is the content 
% Source PDF page 301; original printed page 284
\BookSource{301}{284}
of \BookInline{g_{i}(x)} and \BookInline{h_{i}(x)} is a primitive polynomial. Each of the \BookInline{h_{i}(x)} is irreducible, because if it could be factored, none of the factors could lie in \BookInline{D,\allowbreak } hence all would have positive degree leading to a corresponding factorization of \BookInline{g_{i}(x),\allowbreak } and then to a factorization of \BookInline{f(x)} with more than \BookInline{r} factors of positive degree, contradicting our choice of \BookInline{r.} Thus we now have

\end{BookBody}

\BookNarrativetrue

\begin{BookBody}

\begin{center}\begin{adjustbox}{max width=\linewidth}\BookInline{\displaystyle f(x)= c_{1}h_{1}(x)c_{2}h_{2}(x)\cdots c_{r}h_{r}(x)}\end{adjustbox}\end{center}

\noindent where the \BookInline{h_{i}(x)} are irreducible in \BookInline{D[x].} If we now factor the \BookInline{c_{i}} into irreducibles in \BookInline{D,\allowbreak } we obtain a factorization of \BookInline{f(x)} into a product of irreducibles in \BookInline{D[x].}

The factorization of \BookInline{f(x)\in D[x],\allowbreak } where \BookInline{f(x)} has degree 0, is unique since \BookInline{D} is a UFD; see the comment following Lemma 34.28. If \BookInline{f(x)} has degree greater than 0, we can view any factorization of \BookInline{f(x)} into irreducibles in \BookInline{D[x]} as a factorization in \BookInline{F[x]} into units (that is, the factors in \BookInline{D)} and irreducible polynomials in \BookInline{F[x]} by Lemma 34.28. By Theorem 28.21, these polynomials are unique, except for possible constant factors in \BookInline{F.} But as an irreducible in \BookInline{D[x],\allowbreak } each polynomial of degree \BookInline{>0} appearing in the factorization of \BookInline{f(x)} in \BookInline{D[x]} is primitive. By the uniqueness part of Lemma 34.24, this shows that these polynomials are unique in \BookInline{D[x]} up to unit factors, that is, associates. The product of the irreducibles in \BookInline{D} in the factorization of \BookInline{f(x)} is the content of \BookInline{f(x),\allowbreak } which is again unique up to a unit factor by Lemma 34.24. Thus all irreducibles in \BookInline{D[x]} appearing in the factorization are unique up to order and associates\BookInline{.}\BookEnd{\blacklozenge}

\BookLabel{34.31 Corollary} If \BookInline{F} is a field and \BookInline{x_{1},\allowbreak \cdots,\allowbreak x_{n}} are indeterminates, then \BookInline{F[x_{1},\allowbreak \cdots,\allowbreak x_{n}]} is a UFD.

\BookLabel{\textit{Proof}} By Theorem \BookInline{28.21,\allowbreak F[x_{1}]} is a UFD. By Theorem 34.30, so is \BookInline{(F[x_{1}])[x_{2}]= F[x_{1},\allowbreak x_{2}].} Continuing in this procedure, we see (by induction) that \BookInline{F[x_{1},\allowbreak \cdots,\allowbreak x_{n}]} is a UFD\BookInline{.}\BookEnd{\blacklozenge}

We have seen that a PID is a UFD. Corollary 34.31 makes it easy for us to give an example that shows that\textit{ not every} UFD\textit{ is} \BookInline{a} PID.

\BookLabel{34.32 Example} Let \BookInline{F} be a field and let \BookInline{x} and \BookInline{y} be indeterminates. Then \BookInline{F[x,\allowbreak y]} is a UFD by Corollary 34.30. Consider the set \BookInline{N} of all polynomials in \BookInline{x} and \BookInline{y} in \BookInline{F[x,\allowbreak y]} having constant term 0. Then \BookInline{N} is an ideal, but not a principal ideal. Thus \BookInline{F[x,\allowbreak y]} is not a PID\BookInline{.}\BookEnd{\blacktriangle}

Another example of a UFD that is not a PID is \BookInline{\mathbb{Z}[x],\allowbreak } as shown in Exercise 12, Section 35.

\end{BookBody}

\BookNarrativefalse

\begin{BookBody}

\BookExerciseTitle{34}

\BookTextHeading{\textbf{Computations}}

\noindent In Exercises 1 through 8, determine whether the element is an irreducible of the indicated domain.

\begin{center}\begin{adjustbox}{max width=\linewidth}\BookInline{\displaystyle 2.- 17\;\text{in}\;\mathbb{Z}}\end{adjustbox}\end{center}

\BookExercise{1.}{5 in \BookInline{\mathbb{Z}}}

\BookExercise{3.}{14 in \BookInline{\mathbb{Z}}}

\BookExercise{4.}{\BookInline{2x- 3} in \BookInline{\mathbb{Z}[x]}}

\BookExercise{5.}{\BookInline{2x- 10} in \BookInline{\mathbb{Z}[x]}}

\BookExercise{6.}{\BookInline{2x- 3} in \BookInline{\mathbb{Q}[x]}}

\BookExercise{7.}{\BookInline{2x- 10} in \BookInline{\mathbb{Q}[x]}}

\BookExercise{8.}{\BookInline{2x- 10} in \BookInline{\mathbb{Z}_{11}[x]}}

\BookExercise{9.}{If possible, give four different associates of \BookInline{2x- 7} viewed as an element of \BookInline{\mathbb{Z}[x];} of \BookInline{\mathbb{Q}[x];} of \BookInline{\mathbb{Z}_{11}[x].}}

\BookExercise{10.}{Factor the polynomial \BookInline{4x^{2}- 4x+ 8} into a product of irreducibles viewing it as an element of the integral domain \BookInline{\mathbb{Z}[x];} of the integral domain \BookInline{\mathbb{Q}[x];} of the integral domain \BookInline{\mathbb{Z}_{11}[x].}}

\noindent In Exercises 11 through 13, find all gcd’s of the given elements of \BookInline{\mathbb{Z}.}

\BookExercise{11.}{234, 3250, 1690}

\BookExercise{12.}{\BookInline{784,\allowbreak - 1960,\allowbreak 448}}

\BookExercise{13.}{2178, 396, 792, 594}

\end{BookBody}

% Source PDF page 302; original printed page 285
\BookSource{302}{285}
\BookNarrativefalse

\begin{BookBody}

\noindent In Exercises 14 through 17, express the given polynomial as the product of its content with a primitive polynomial in the indicated UFD.

\BookExercise{14.}{\BookInline{18x^{2}- 12x+ 48} in \BookInline{\mathbb{Z}[x]}}

\BookExercise{15.}{\BookInline{18x^{2}- 12x+ 48} in \BookInline{\mathbb{Q}[x]}}

\BookExercise{16.}{\BookInline{2x^{2}- 3x+ 6} in \BookInline{\mathbb{Z}[x]}}

\BookExercise{17.}{\BookInline{2x^{2}- 3x+ 6} in \BookInline{\mathbb{Z}_{7}[x]}}

\BookTextHeading{\textbf{Concepts}}

\noindent In Exercises 18 through 20, correct the definition of the italicized term without reference to the text, if correction is needed, so that it is in a form acceptable for publication.

\BookExercise{18.}{Two elements \BookInline{a} and \BookInline{b} in an integral domain \BookInline{D} are\textit{ associates} in \BookInline{D} if and only if their quotient \BookInline{a/b} in \BookInline{D} is a unit.}

\BookExercise{19.}{An element of an integral domain \BookInline{D} is an\textit{ irreducible} of \BookInline{D} if and only if it cannot be factored into a product of two elements of \BookInline{D.}}

\BookExercise{20.}{An element of an integral domain \BookInline{D} is a\textit{ prime} of \BookInline{D} if and only if it cannot be factored into a product of two smaller elements of \BookInline{D.}}

\BookExercise{21.}{Determine whether each of the following is true or false.}

\BookExercise{a.}{Every field is a UFD.}

\BookExercise{b.}{Every field is a PID.}

\BookExercise{c.}{Every PID is a UFD.}

\BookExercise{d.}{Every UFD is a PID.}

\BookExercise{e.}{\textbf{e}\BookInline{.\mathbb{Z}[x]} is a UFD.}

\BookExercise{f.}{Any two irreducibles in any UFD are associates.}

\BookExercise{g.}{If \BookInline{D} is a PID, then \BookInline{D[x]} is a PID.}

\BookExercise{h.}{If \BookInline{D} is a UFD, then \BookInline{D[x]} is a UFD.}

\BookExercise{i.}{In any UFD, if \BookInline{p|a} for an irreducible \BookInline{p,\allowbreak } then \BookInline{p} itself appears in every factorization of \BookInline{a.}}

\BookExercise{j.}{A UFD has no divisors of 0.}

\BookExercise{22.}{Let \BookInline{D} be a UFD. Describe the irreducibles in \BookInline{D[x]} in terms of the irreducibles in \BookInline{D} and the irreducibles in \BookInline{F[x],\allowbreak } where \BookInline{F} is a field of quotients of \BookInline{D.}}

\BookExercise{23.}{Lemma 34.28 states that if \BookInline{D} is a UFD with a field of quotients \BookInline{F,\allowbreak } then a nonconstant irreducible \BookInline{f(x)} of \BookInline{D[x]} is also an irreducible of \BookInline{F[x].} Show by an example that a \BookInline{g(x)\in D[x]} that is an irreducible of \BookInline{F[x]} need not be an irreducible of \BookInline{D[x].}}

\BookExercise{24.}{All our work in this section was restricted to integral domains. Taking the same definition in this section but for a commutative ring with unity, consider factorizations into irreducibles in \BookInline{\mathbb{Z}\times \mathbb{Z}.} What can happen? Consider in particular (1, 0).}

\BookTextHeading{\textbf{Theory}}

\BookExercise{25.}{Prove that if \BookInline{p} is a prime in an integral domain \BookInline{D,\allowbreak } then \BookInline{p} is an irreducible.}

\BookExercise{26.}{Prove that if \BookInline{p} is an irreducible in a UFD, then \BookInline{p} is a prime.}

\BookExercise{27.}{For a commutative ring \BookInline{R} with unity show that the relation \BookInline{a\sim b} if \BookInline{a} is an associate of \BookInline{b(}that is, if \BookInline{a= bu} for \BookInline{u} a unit in \BookInline{R)} is an equivalence relation on \BookInline{R.}}

\BookExercise{28.}{Let \BookInline{D} be an integral domain. Exercise 39, Section 22 showed that \BookInline{\langle U,\allowbreak \cdot \rangle } is a group where \BookInline{U} is the set of units of \BookInline{D.} Show that the set \BookInline{D^{\ast }- U} of nonunits of \BookInline{D} excluding 0 is closed under multiplication. Is this set a group under the multiplication of \BookInline{D}?}

\BookExercise{29.}{Let \BookInline{D} be a UFD. Show that a nonconstant divisor of a primitive polynomial in \BookInline{D[x]} is again a primitive polynomial.}

\BookExercise{30.}{Show that in a PID, every proper ideal is contained in a maximal ideal. [\textit{Hint:} Use Lemma 34.10.]}

\BookExercise{31.}{Factor \BookInline{x^{3}- y^{3}} into irreducibles in \BookInline{\mathbb{Q}[x,\allowbreak y]} and prove that each of the factors is irreducible.}

\noindent There are several other concepts often considered that are similar in character to the ascending chain condition on ideals in a ring. The following three exercises concern some of these concepts.

\end{BookBody}

% Continuation on source PDF page 303, preceding Section 35
\BookNarrativefalse

\begin{BookBody}

\BookExercise{32.}{Let \BookInline{R} be any ring. The\textbf{ ascending chain condition} (ACC)\textbf{ for ideals} holds in \BookInline{R} if every strictly increasing sequence \BookInline{N_{1}\subset N_{2}\subset N_{3}\subset \cdots } of ideals in \BookInline{R} is of finite length. The\textbf{ maximum condition} (MC)\textbf{ for ideals} holds in \BookInline{R} if every nonempty set \BookInline{S} of ideals in \BookInline{R} contains an ideal not properly contained in any other ideal of the set \BookInline{S.} The\textbf{ finite basis condition} (FBC)\textbf{ for ideals} holds in \BookInline{R} if for each ideal \BookInline{N} in \BookInline{R,\allowbreak } there is a finite set \BookInline{B_{N}= \{b_{1},\allowbreak \cdots,\allowbreak b_{n}\}\subseteq N} such that \BookInline{N} is the intersection of all ideals of \BookInline{R} containing \BookInline{B_{N}.} The set \BookInline{B_{N}} is a\textbf{ finite} \textbf{generating set for} \BookInline{N.}}

Show that for every ring \BookInline{R,\allowbreak } the conditions ACC, MC, and FBC are equivalent.

\BookExercise{33.}{Let \BookInline{R} be any ring. The\textbf{ descending chain condition} (DCC)\textbf{ for ideals} holds in \BookInline{R} if every strictly decreasing sequence \BookInline{N_{1}\supset N_{2}\supset N_{3}\supset \cdots } of ideals in \BookInline{R} is of finite length. The\textbf{ minimum condition} (mC)\textbf{ for ideals} holds in \BookInline{R} if given any set \BookInline{S} of ideals of \BookInline{R,\allowbreak } there is an ideal of \BookInline{S} that does not properly contain any other ideal in the set \BookInline{S.}}

Show that for every ring, the conditions DCC and mC are equivalent.

\BookExercise{34.}{Give an example of a ring in which ACC holds but DCC does not hold. (See Exercises 32 and 33.)}

\end{BookBody}
\ifdefined\BookMaster\let\BookFinish\relax\else\def\BookFinish{\end{document}}\fi\BookFinish
